\section{Correctness}
\label{sec:verification}

In the previous sections, we discussed removing compression as a design choice with engineering consequences. We now show why that
change matters beyond engineering and show that it turns the decryption-failure probability from a heuristic estimate into a
formally verified bound.

\paragraph{Why compression is the obstruction.}
In the noise expression that Barbosa et al.~\cite{BarbosaKLSS25} derive for ML-KEM,
$\tilde{n} = \langle e,r\rangle - \langle s,e_1\rangle - \langle s,c_u\rangle + e_2 + c_v$,
the rounding errors $c_u$ and $c_v$ are the only source of difficulty. First, because the correctness adversary holds the secret key, 
no direct MLWE reduction is available which forces a worst-case elimination of $c_v$ and a split into two sub-bounds. Second, $c_u$'s own rounding-error 
distribution is not symmetric and so fails the ``good'' property their framework requires. This forces an index-dependent, 256-term sum rather than a 
single closed-form union bound. In the case of FrodoKEM, since it admits no ciphertext compression, the noise term only involves
good distributions, which gives a clean union bound with no heuristic assumption. The cost of compression is concrete and quantitative. For ML-KEM-768, 
the fully provable bound (which pays for compression) is $2^{-80}$, while the heuristic bound that just assumes rounding-of-uniform is 
$2^{-164}$, a gap of roughly 84 bits.

\paragraph{Consequence for \textsc{Kaiburr}}
We adapt the EasyCrypt framework of~\cite{BarbosaKLSS25} and instantiate it for \textsc{Kaiburr} which has no ciphertext compression like FrodoKEM. 
We reuse their theory for
combining distributions over a ring and their proofs that the relevant properties are preserved. We also use their reduction to a 
single coefficient and their OCaml implementation for computing a sound upper bound on the resulting tail probability. 
\textsc{Kaiburr} retains \textsc{Kyber} as specified in
Sections~\ref{sec:preliminaries}--\ref{sec:implementation} in every respect the analysis depends on. The only
property of the noise family $f(t)$ the proof uses is that it is ``good'' in their sense: non-trivial and symmetric around 0 
(see Section~\ref{sec:distribution}).

Since the rounding terms vanish, the noise
collapses to $\tilde{n} = \langle e,r\rangle - \langle s,e_1\rangle + e_2$. As a result, every coefficient is identically distributed independent of its
index, and the union bound becomes a single factor of $256$ rather than a sum of $256$ distinct terms. Moreover, no additional
MLWE term is needed. We note the resulting expression is in fact simpler than the one obtained for FrodoKEM, since
$q=3329$ is prime, so $\mathbb{Z}_q$ is a field and goodness is preserved through inner products with only $\oplus$
required, whereas FrodoKEM's power-of-two modulus forces $\ominus$ to be retained.

\paragraph{Computing the bound.}
The tail probability is computed outside EasyCrypt by OCaml code that mirrors the
EasyCrypt expression. We report the results for kaiburr-1792 ($k=7$), kaiburr-4608 ($k=18$), and kaiburr-6144 ($k=24$) in Table~\ref{tab:bounds},
alongside the failure probability obtained for each parameter set from the \textsc{Kyber} team's estimator scripts
(Table~\ref{tab:kaiburr-params}).

\begin{table}[h]
  \centering
  \caption{Decryption-failure probability for each \textsc{Kaiburr} parameter set: the formally verified bound
  obtained in this work, and the estimate from the \textsc{Kyber} team's scripts (Table~\ref{tab:kaiburr-params}).}
  \label{tab:bounds}
  \begin{tabular}{lccc}
    \toprule
    Scheme & $k$ & Verified bound (this work) & \textsc{Kyber}-script estimate \\
    \midrule
    kaiburr-1792 & 7  & $2^{-133.3}$ & $2^{-133.7}$ \\
    kaiburr-4608 & 18 & $2^{-133.8}$ & $2^{-134.1}$ \\
    kaiburr-6144 & 24 & $2^{-139.6}$ & $2^{-139.9}$ \\
    \bottomrule
  \end{tabular}
\end{table}

To put the numbers in context: for ML-KEM-768, Barbosa et al.~\cite{BarbosaKLSS25}
compute a heuristic bound of $2^{-164}$, while the best bound they can
justify against the definition of correctness is only
$2^{-80}$, a gap of roughly 84 bits. For \textsc{Kaiburr} there is no such gap to close, since removing compression
leaves no rounding term for a heuristic assumption to apply to. Table~\ref{tab:bounds} shows our provable bound
closely tracks the \textsc{Kyber} team's own estimate for every parameter set. This shows that the bound is provable as well
as tight in practice.

The price \textsc{Kaiburr} pays by removing compression is ciphertext size. However, this additional size is modest in the context of the high-security 
parameter sets considered in this work; it also remains within our fixed communication budget.